\section{Testing Defense Transfer Boundaries}\label{sec:controlled-placement}\label{sec:tau-cover-bridge}
Four experiments instantiate the boundaries of \S\ref{sec:discussion} using saved LLM responses and a retail agent workload: the emitter (Exp~1), executed state (Exp~2), calibration population (Exp~3) and attacker observation (Exp~4). Exps~1--2 reproduce known mechanisms to measure their streaming-cost and state/link consequences; Exps~3--4 measure agent-specific operating points. Table~\ref{tab:r7-cohorts} separates populations; Appendix~\ref{app:experiments} indexes all studies.

\paragraph{Shared measurement conventions}
A byte-transparent relay captures post-handshake TLS application records on loopback; TLS~1.3 session tickets are disabled. Standard-prompt results use the arm called \texttt{baseline} in the artifact; \emph{undefended} denotes transport without a program defense. Balanced accuracy (BA) averages the two class recalls. We report the maximum over each specified attack suite and recompute that maximum in each of 2,000 common bootstrap draws, grouping tasks for reference-program tests and customers for real-agent tests. The original protocol's engineering screen requires an upper 95\% BA bound below 0.60 and maximum AUROC below 0.65 for both targets. No defended configuration in Exps~3--4 passes it, so we report paired reductions against named suites rather than absolute protection (\S\ref{sec:validity}).

\paragraph{Attack suites}
Because BA is a suite maximum, every number is indexed by its suite. The \emph{frozen suite} (Exp~3 frozen) contains clean-trained $D_0$, defended-refit $D_2$ and training-derived decoder/count rules over lengths and order. The \emph{11-attack union} (Exp~3 cross-fit and localization) adds learners including a 1D-CNN and a k-FP-style forest-leaf model. The \emph{endpoint suite} (Exp~3 replications and transfer) adds a total-exchange-count rule to six attacks inherited from the frozen protocol (majority, clean- and defended-trained learners, decoded-tool and target-call-count rules). In Exp~4, the \emph{common seven} are learners shared with the offline WF adaptations, and the \emph{full union} (17--22 attacks) adds timing features and learned flush removal.
\begin{table}[t]\centering\footnotesize
\caption{Cohort, observation and cost denominators for Exps~3--4. Standard-prompt pairs: cancel/return BA (\%). Byte ratios price joint shaping against each study's own undefended serialization.}\label{tab:r7-cohorts}
\begin{tabularx}{\columnwidth}{@{}lYrr@{}}\toprule
Study & Observation / attack suite & Undef.\ BA & Joint B$\times$\\\midrule
Exp~3 frozen & Replay lengths/order; frozen suite & 97.2/94.0 & 15.39\\
Exp~3 cross-fit & Analytic lengths/order; 11-attack union & 91.7/85.5 & 19.76\\
Exp~4 & New RPC envelope; lengths/order/time; full union & 91.7/86.4 & 15.11\\\bottomrule
\end{tabularx}
\par\smallskip\raggedright\footnotesize Frozen-split and Exp~4 rows: 276 test episodes, 24 customers, 5/7 positive customers. Exp~3 cross-fit: 690 episodes, 53 customers, 14/18 positive customers. Each prompt has half the episodes. Cross-fit fold calibration yields larger templates ($5.04\times$ calls); frozen-split and Exp~4 rows use $3.83\times$. Repeats measure execution variation within these customer groups.
\end{table}


\subsection{Exp 1: bind protection and cost to the emitter}\label{sec:exp-emitter}
Buffered tests use brief/expanded saved responses for 24 topics (12 training, 12 test). Padding only the upstream link leaves verbosity 100\% separable at a plaintext-aware terminal emitter; padding both reduces BA to 50\%, reproducing NetEcho's mechanism~\cite{netecho2025}. A training-fitted buffered cap achieves 50\% at 2.18$\times$ wire bytes under TLS~1.2 (2.19$\times$ under TLS~1.3).

Streaming the same 48 texts as per-token server-sent events changes the cost contract~\cite{weiss2024prompt,whisperleak2025}. Per-event padding retains token count and 100\% verbosity inference. Eight-token batching cuts bytes to 0.21$\times$ and clean-trained BA to 50\%, yet defended-data refitting restores 100\%. Constant shape reaches 50\% at 2.42$\times$ bytes and 208 versus 105 token intervals; two overflowing test outputs remain identifiable (Table~\ref{tab:transfer-streaming}).
\begin{table}[t]\centering\footnotesize\setlength{\tabcolsep}{1.9pt}
\caption{Exp 1: the same 48 saved responses emitted by a buffered or a token-streaming frontend (TLS~1.3 loopback captures). Bytes are relative to the unpadded emitter of the same kind. TTFB/Done: first and last delivered byte in token intervals (accounting proxy). Accuracies on 24 held-out outputs (12 topics); ID: closed-world top-1 identification among 48 responses (chance 2.1\%).}\label{tab:transfer-streaming}
\begin{tabular}{@{}lrrrrrrr@{}}\toprule
& & & & \multicolumn{2}{c}{Verbosity} & \multicolumn{2}{c}{Prefix}\\\cmidrule(lr){5-6}\cmidrule(l){7-8}
Emitter / framing & Bytes$\times$ & TTFB & Done & $D_0$ & $D_2$ & $D_2$ & ID\\\midrule
Buffered, unpadded & 1.00 & 105 & 105 & 100.0 & 100.0 & 91.7 & 87.5\\
Buffered, training cap & 2.19 & 105 & 105 & 50.0 & 50.0 & 50.0 & 2.1\\
Stream, per token & 1.00 & 1 & 105 & 100.0 & 100.0 & 87.5 & 100.0\\
\quad+ random event pad & 1.29 & 1 & 105 & 100.0 & 100.0 & 87.9 & 62.3\\
\quad+ constant event size & 1.23 & 1 & 105 & 100.0 & 100.0 & 87.5 & 58.3\\
8-token batches, const. & 0.21 & 8 & 105 & 50.0 & 100.0 & 87.5 & 20.8\\
Const.\ events + count cap & 2.42 & 1 & 208 & 50.0 & 50.0 & 50.0 & 6.3\\\bottomrule
\end{tabular}
\par\smallskip\raggedright\footnotesize $D_2$ is the best of total-byte threshold, record-count threshold and a logistic sequence model, each fitted on defended training outputs. Random padding: 10 socket-measured draws plus an independent profiling run. TLS~1.2 gives identical accuracies and byte ratios within 0.03. Verbosity contrasts brief/expanded outputs; Prefix contrasts the first two/three sentences of each expanded output. Tokens use \texttt{o200k\_base}. Training count caps are 208/88 events, with 2/24 and 5/24 test overflows, respectively.
\end{table}

\textbf{Finding.} The same saved content yields different protection and delay under buffering and streaming. A buffered cap's byte ratio cannot price a streaming deployment, whose first-output delay and emission schedule are different outcomes.

\subsection{Exp 2: compose state preservation with link protection}\label{sec:exp-state}
\paragraph{Execution boundary}
Among 115 public retail tasks, 97 have an initially pending order for the task user. Three order-selection seeds yield 291 cover placements before the reference program. Same-item modification changes status and appends a refund in 173 placements despite unchanged item and price; the other 118 calls fail the variant-availability check. Idempotence, $f(f(s,a),a)=f(s,a)$, therefore differs from state preservation, $f(s,a)=s$. Cancelling another user's order changes the shared database. The shadow executor isolates cover effects on a separate state and suppresses its responses; a live-execution control exposes the resulting side effects.

The dispatch observer sees ordered tool names and response lengths. Policies add cancellation, pad its count, interleave donor programs or equalize all tool counts; donors/templates use the known 115-program population. Task-grouped five-fold attacks fit clean/defended gradient boosting and logistic regression on counts, bigrams and lengths, with a count-rule baseline and three seeds (18/32 positive cancellation/return tasks).
\begin{table}[t]\centering\footnotesize\setlength{\tabcolsep}{3pt}
\caption{Exp 2 dispatch results: all 115 reference programs, three seeds, known-population defense calibration; task-folded attackers. Maximum BA (cancel/return); extra calls per real call; database equality with the unprotected run.}\label{tab:r7-e6}
\begin{tabularx}{\columnwidth}{@{}Yrrr@{}}\toprule
Policy & BA (\%) & Calls+ & DB eq.\\\midrule
No cover & 100.0/100.0 & 0.00 & 100\%\\
Cancel, $p=.5$ & 86.1/100.0 & .10 & 100\%\\
Cancel, $p=1$, shadow & 100.0/100.0 & .20 & 100\%\\
Cancel, $p=1$, live & 100.0/100.0 & .20 & 0\%\\
Cancel count $\to2$ & 76.9/100.0 & .35 & 100\%\\
One decoy program & 91.9/85.9 & 1.01 & 100\%\\
Two decoy programs & 83.5/78.1 & 1.98 & 100\%\\
Shape equalization & 50.0/54.5 & 6.71 & 100\%\\\bottomrule
\end{tabularx}\end{table}

One inserted cancellation drives clean-trained BA to chance, while refitting restores 100\%. Count equalization removes that rule's signal, but correlated tools retain 76.9\% BA. Donor programs, inspired by WRA's discarded retrieval cover~\cite{jeong2026wra}, leave 83.5--91.9\%. Shape equalization reaches 50.0/54.5\% at 6.7 extra calls per real call (Table~\ref{tab:r7-e6}). Shadow checks match the original live state in every case; live cancellation changes it in every case.

\paragraph{Final network view}
Four cover policies and three framing choices produce 4,140 TLS~1.2 sessions. The network observer receives direction-tagged record lengths and order. A training-fold profiler maps request sizes to tool types before applying dispatch attacks: request-only identification reaches 89.9/91.4\% (582 calls; majority 29.4\%). Per-call buckets leave return BA at 73.9\%; joint shape cover and constant framing reduce cancellation/return to 54.4/50.0\%, at 26.3$\times$ bytes and 7.7$\times$ calls.

\textbf{Finding.} Call-level cover and per-call framing address complementary observables. Full constant-rate transport can instead hide packet counts from a passive link observer without executing dummy tools~\cite{rfc9347}; Exp~4 measures that alternative. At a dispatch view after decapsulation, cover requires hidden markers and routing to the trusted isolated executor (\S\ref{sec:deployment}). The intervention must match the receiver and the protected state.

\subsection{Exp 3: real-agent signatures and calibrated shaping}\label{sec:real-agent}\label{sec:r5}
\paragraph{Task and collection}
The targets are the reference program's requested operations, \texttt{cancel\_pending\_order} and \texttt{return\_delivered\_order\_items}. This operational target tests whether encrypted metadata reveals which service action a task requests. GPT-6.1 Sol/low plans from public retail policies, schemas and actual sandbox observations; a separate context simulates the user. Two prompts and three repeats on 115 tasks yield 690 runs from 53 customers. The rules prompt adds a customer-details read before order access and a read before each write.

\paragraph{Frozen calibration and behavioral control}
All 414 training runs (69 tasks) precede freezing shared cover templates, donors and framing caps. The 276 test runs cover 46 tasks from 24 disjoint customers. Prefix-only scheduling inserts training-derived cover, flushing unmet counts at episode end. Four program conditions per prompt give 2,760 local TLS replays with lengths/order attacks; every failure and handoff remains included. Real-call pass-through plus isolated cover preserves the logical response/state projection (Appendix~\ref{app:equivalence}); reusing trajectories pairs the network interventions on identical model decisions.

Rules raise read-order adherence from 86.2\% to 97.1\%, while undefended category inference remains above 93\%. Native rewards are 110/138 and 105/138 (paired difference $-3.6$ points, 95\% interval $[-9.5,1.4]$). Thus compliance and metadata privacy separate in the same task population.
\begin{table*}[t]\centering\small
\caption{Exp 3, frozen calibration: operation-category inference on real-agent trajectories. Testing covers 46 tasks from 24 held-out customers, two prompts and three repeats (276 episodes). Network BA is the maximum over the frozen attack suite, with customer-cluster 95\% intervals.}\label{tab:r4-primary}
\begin{tabular}{@{}lrrrrr@{}}\toprule
Condition & Cancel BA [95\% CI] & Return BA [95\% CI] & Original reward (\%) & Bytes$\times$ & Calls$\times$\\\midrule
Undefended & 97.2 [90.0, 100.0] & 94.0 [85.8, 100.0] & 79.7 & 1.00 & 1.00\\
Rules prompt & 97.2 [90.0, 100.0] & 93.5 [85.0, 99.5] & 76.1 & 1.00 & 1.00\\
Padding only & 66.1 [50.0, 83.9] & 55.5 [50.0, 73.0] & -- & 4.02 & 1.00\\
Shape only & 60.8 [55.4, 71.3] & 51.9 [50.5, 57.8] & -- & 4.16 & 3.83\\
Joint program & 50.0 [50.0, 60.6] & 50.0 [50.0, 65.5] & -- & 15.39 & 3.83\\
Rules + joint & 50.0 [50.0, 65.1] & 60.6 [50.0, 73.3] & -- & 15.72 & 3.88\\
\bottomrule\end{tabular}
\par\smallskip\raggedright \textbf{Paired joint-minus-undefended BA differences (percentage points):} standard-prompt cancellation $-47.2$ [$-50.0$, $-33.3$], return $-44.0$ [$-48.9$, $-27.8$]. These post hoc contrasts use the same customer draws and reselect the attack-suite maximum in each condition. Joint = online shadow cover plus constant framing; padding/shape ablations use standard-prompt trajectories. Ratios use undefended transport within each prompt. Shadow isolation preserves the logical planner transcript; the four program conditions replay actual LLM trajectories over local TLS~1.2. Reward is reported only for original model executions; dashes denote no new live defended-planner reward measurement (\S\ref{sec:validity}).
\end{table*}

Joint cover and framing reach 50.0/50.0\% point BA, compared with 97.2/94.0\% undefended, a paired decrease of 47.2/44.0 points (Table~\ref{tab:r4-primary}). Their upper bounds (60.6/65.5\%) fail the screen, and the rules-arm joint return point is 60.6\%, identifying prompt/cohort sensitivity.

\paragraph{Customer-grouped exploratory comparison}
Five-fold cross-fitting refits templates, donors, caps and all attackers on other customers, pooling the full known cohort (14/18 positive customers). The analytic TLS model matches all 2,760 saved sessions and 120 additional captures. Under the 11-attack union, joint constant shaping reduces BA by 31.2/30.0 points (paired 95\% intervals [15.3,41.0]/[14.8,39.2]); buckets reduce it by 34.6/32.8 points, at lower byte cost (Table~\ref{tab:r5}).
\begin{table}[t]\centering\footnotesize\setlength{\tabcolsep}{2pt}
\caption{Exp 3: exploratory cross-fit comparison under the 11-attack union: all 690 episodes and 53 customers. Maximum network BA with customer-bootstrap 95\% intervals. Standard prompt unless marked Rules; byte/call ratios relative to undefended transport.}\label{tab:r5}
\begin{tabular}{@{}lrrr@{}}\toprule
Configuration & Cancel BA (\%) & Return BA (\%) & B$\times$/C$\times$\\\midrule
Undefended & \mbox{91.7 [82.1, 100.0]} & \mbox{85.5 [76.7, 93.4]} & 1.0/1.0\\
Joint constant & \mbox{60.5 [52.9, 72.4]} & \mbox{55.5 [51.8, 64.2]} & 19.8/5.0\\
Joint buckets & \mbox{57.1 [50.0, 71.5]} & \mbox{52.7 [50.0, 61.7]} & 12.4/5.0\\
Rules + buckets & \mbox{56.9 [50.2, 71.2]} & \mbox{52.7 [50.0, 65.0]} & 12.4/5.0\\
\bottomrule\end{tabular}
\par\smallskip\raggedright\footnotesize Joint constant is the frozen study's shadow-count plus constant-framing design, refit within each fold. Joint buckets changes framing to 1024-byte buckets. Every bootstrap draw reselects the complete maximum. Intervals condition on fitted predictions.
\end{table}


\paragraph{Locate the observable operation signature}
An oracle selects write exchanges using true tool types, then exposes only ciphertext lengths/order. Writes comprise 19.6\% of real calls yet retain 91.7/84.9\% BA; removing them leaves 62.8/59.9\% (Table~\ref{tab:r7-e12}). The target names an operation, so the result locates its metadata signature. It directs protection toward specific exchanges, while residual read inference shows why hiding one named call is insufficient.
\begin{table}[t]\centering\footnotesize\setlength{\tabcolsep}{2pt}
\caption{Exp 3 oracle localization under the 11-attack union. Standard prompt: 345 cross-fitted episodes, 53 customers. Maximum network BA (\%) with customer-bootstrap intervals.}\label{tab:r7-e12}
\begin{tabular}{@{}lrrr@{}}\toprule
Target & Full trace & Writes removed & Writes only\\\midrule
Cancel & \mbox{91.7 [82.1, 100.0]} & \mbox{62.8 [54.6, 75.3]} & \mbox{91.7 [82.1, 100.0]}\\
Return & \mbox{85.5 [76.7, 93.4]} & \mbox{59.9 [52.2, 71.8]} & \mbox{84.9 [78.8, 94.0]}\\
\bottomrule\end{tabular}\end{table}


\paragraph{Model-specific repeated evaluation}
Luna's 150-run extension contributes a same-provider replication; two additional endpoints each contribute 115 standard-prompt runs with 69 training tasks, a defense frozen before test collection, and 46 test tasks (Table~\ref{tab:r11-model-replication}). Joint shaping lowers return BA by 16.0 points for MiniMax-M3 and 19.4 for GLM; cancellation changes by $+2.1$ and $-50.0$ points, respectively. Return-operation inference thus recurs under separate within-model fits, while protection depends on model and target. MiniMax-M3 is a served alias (its operator reports a DeepSeek backend); GLM uses a 16,384-token budget versus 4,096 for MiniMax; neither endpoint exposes an immutable revision.
\begin{table}[t]\centering\small
\caption{Exp 3 model-specific repeated evaluations: network suite-max BA (cancel/return, \%). Each model has 46 test tasks and 24 customers. Endpoint rows use the endpoint suite; Luna retains the frozen suite. Luna uses 29 training tasks; the two endpoint replications use 69. Contrasts are within model; endpoint aliases and inference budgets differ.}\label{tab:r11-model-replication}
\begin{tabular}{@{}lrr@{}}\toprule
Model & Undefended & Joint\\\midrule
Luna & 64.2 / 84.9 & 60.4 / 57.0\\
MiniMax-M3 & 58.3 / 83.5 & 60.4 / 67.5\\
GLM-5.3-flash & 100.0 / 86.6 & 50.0 / 67.3\\
\bottomrule\end{tabular}
\end{table}

\paragraph{Separate calibration source from sampling sensitivity}
Table~\ref{tab:calibration-transfer} swaps source calibration and refits the complete endpoint suite on target training customers. On GLM, the imported template changes cancellation/return BA by $+20.8/-17.3$ points and increases bytes. Customer-paired permutation tests do not resolve either direction (raw $p=0.134/0.174$; Holm-adjusted $p=1.000$ across all eight contrasts). Rebuilding calibration over 100 paired training-customer resamples gives local GLM return BA a 95\% sensitivity range of 50.0--67.9\%, comparable in width to the original 17.3-point contrast. The imported-minus-local return difference has median 0.0 points and spread $[-16.9,18.4]$; with attacker resampling it is $+0.3$ $[-11.2,18.9]$. Thus the original direction is not stable against local training variability.
\begin{table*}[t]\centering\footnotesize
\caption{Calibration transfer and write-only padding on each target model\textquotesingle s 46 test tasks (24 customers). BA is the endpoint-suite maximum; paired $\Delta$ is relative to target-local joint shaping (percentage points). Costs divide summed test bytes/calls by undefended target-model totals.}\label{tab:calibration-transfer}
\begin{tabularx}{\textwidth}{@{}llrrYYrr@{}}\toprule
Target & Calibration / intervention & Cancel BA & Return BA & $\Delta$ cancel [95\% CI] & $\Delta$ return [95\% CI] & Bytes$\times$ & Calls$\times$\\\midrule
MiniMax & MiniMax joint & 60.4 & 67.5 & -- & -- & 20.06 & 3.36\\
MiniMax & GLM joint & 50.8 & 63.0 & $-9.6$ [-21.1, 6.4] & $-4.5$ [-25.0, 12.3] & 17.48 & 3.11\\
MiniMax & Write padding & 50.0 & 61.3 & $-10.4$ [-23.3, 0.0] & $-6.2$ [-20.0, 9.4] & 1.45 & 1.00\\
GLM & GLM joint & 50.0 & 67.3 & -- & -- & 15.95 & 4.24\\
GLM & MiniMax joint & 70.8 & 50.0 & $+20.8$ [-12.6, 35.3] & $-17.3$ [-38.2, -5.3] & 18.65 & 4.66\\
GLM & Write padding & 60.0 & 61.6 & $+10.0$ [-21.4, 34.4] & $-5.7$ [-26.0, 15.8] & 1.51 & 1.00\\
\bottomrule\end{tabularx}
\par\smallskip\raggedright The seven attacks include protected refits, decoded-tool, target-call-count and total-exchange-count rules. Nominal intervals reselect both maxima on 2,000 common customer draws, conditional on the fitted attackers. Each source calibration uses only source training data; target attackers use target training data. Write padding executes no cover calls. These are local TLS replays of fixed model trajectories; no live planner latency or utility is inferred from these costs.
\end{table*}

Write-only padding costs 1.45/1.51$\times$ bytes for MiniMax/GLM but leaves return BA at 61.3/61.6\% under the endpoint suite. A learned length-based write detector has precision 34.9/98.3\% and recall 96.3/89.4\%, distinguishing oracle localization from implementable inference. Localization thus proposes economical interventions; residual signals and detector errors determine whether their protection transfers.

\paragraph{Test live planner utility at its own boundary}
A paired 50-task GLM study runs 100 fresh planner episodes with real isolated prefix-cover execution and emulated response deadlines. Baseline and prefix-shadow succeed on 37/50 and 33/50 tasks (difference -8.0 points, customer interval [-23.9,6.5]). Withheld executed-tool responses trigger subsequent planner actions. This measures the prefix scheduler's utility under injected delays; terminal flush and TLS framing belong to a different mechanism. The supplement retains every request, fault and outcome.

\textbf{Finding.} Read-order compliance and operation privacy separate. Write-local shaping trades lower cost against residual return inference. Calibration comparisons must also account for local recalibration and attacker-fit variability before attributing a change to model source. Table~\ref{tab:r19-decisions} makes these matched comparisons a deployment requirement.

\subsection{Exp 4: timing, latency and matched transport controls}\label{sec:live-timing}\label{sec:r7-main}
All 690 saved action trajectories execute through a live local tool server under six conditions, producing 4,140 new TLS sessions. Each real response gates the next action; real and shadow calls have separate states. The original 414/276 customer-disjoint split is retained, with attacks fitted and frozen before test capture. Lengths, directions and arrival times measure the tool RPC chain.

A complete tunnel sends 4,096-byte cells in both directions at 2 or 5\,ms target intervals; dummy cells execute no tools. A training-sized 133-slot horizon continues after the last response. Serial and two-in-flight shadow schedules use the frozen prefix-only count policy with constant per-call framing. The full union covers size/order, timing, combined features and learned flush removal; the common seven, including a 1D-CNN and forest-leaf neighbors, permit a matched comparison with offline WF adaptations.
\begin{table*}[t]\centering\footnotesize
\caption{Timing-aware Exp 4 comparison: standard prompt, 138 held-out episodes, 24 customers (5 cancel-positive, 7 return-positive). Each BA pair is cancel/return (\%). The common seven models permit a matched-suite comparison; the full union also retains all earlier attacks.}\label{tab:r7-controls}
\begin{tabular}{@{}llrrrrrr@{}}\toprule
View & Condition & Common-seven BA & Full-union BA & Upper BA & Bytes$\times$ & Calls$\times$ & RPC p95 (ms)\\\midrule
M & Undefended & 84.3/76.9 & 91.7/86.4 & 100.0/97.2 & 1.00 & 1.00 & 15.4\\
M & Per-call constant & 53.2/59.1 & 68.6/62.7 & 92.8/79.5 & 3.92 & 1.00 & 16.6\\
M & Serial shadow & 50.0/52.9 & 61.4/53.4 & 71.4/65.8 & 15.11 & 3.83 & 26.2\\
M & Two-in-flight shadow & 50.0/57.3 & 51.7/57.7 & 69.1/69.7 & 15.11 & 3.83 & 20.8\\
M & Full tunnel, 2 ms & 59.9/55.2 & 62.5/55.2 & 70.2/62.0 & 125.71 & 1.00 & 56.4\\
M & Full tunnel, 5 ms & 54.6/53.7 & 54.6/57.3 & 72.9/64.2 & 125.71 & 1.00 & 112.3\\
\midrule
S & Record buckets & 53.1/69.8 & 53.1/69.8 & 78.5/86.9 & 2.25 & 1.00 & --\\
S & Tamaraw 40/12 ms & 55.3/55.3 & 55.3/55.3 & 68.7/64.4 & 12.15 & 1.00 & --\\
S & Tamaraw 4/1.2 ms & 57.8/51.8 & 57.8/51.8 & 71.8/61.8 & 11.99 & 1.00 & --\\
S & FRONT & 59.0/69.3 & 59.0/69.3 & 69.9/81.8 & 2.51 & 1.00 & --\\
S & WTF-PAD adaptation & 51.9/62.7 & 51.9/62.7 & 64.1/79.6 & 1.88 & 1.00 & --\\
S & WT random-donor & 50.0/56.7 & 50.0/56.7 & 67.0/67.8 & 2.04 & 1.00 & --\\
S & FRONT, short calibration & 52.4/56.0 & 52.4/56.0 & 61.3/71.1 & 2.33 & 1.00 & --\\
\bottomrule\end{tabular}
\par\smallskip\raggedright M: measured local tool RPC on saved actions; S: offline traffic transformation. RPC p95 ends at the last real response, excluding LLM generation. S-row simulated spans are not RPC latency. Upper BA is the full-union 95\% bound; every condition exceeds 60\%. M unions contain 17--22 attacks; S unions contain the common seven. Ratios use prompt-specific undefended captures.
\end{table*}


\paragraph{Timing reveals distinguishable cover}
For standard-prompt serial shadow, adding timing/flush attacks to the size/order attacks increases cancellation BA from 50.0\% to 61.4\% (paired increase 11.4 points, 95\% interval [3.7,19.6]); return changes from 50.7\% to 53.4\%. Flush-phase detection reaches 90.4/92.3\% for serial/parallel schedules over 15,278 records from 276 test episodes. A separate cover-versus-real request classifier reaches 82.5/96.1\% on 7,639 requests.

We then compose the detector with intent inference: a training-fitted detector strips the predicted suffix before classification. Table~\ref{tab:r8-flush} shows the recovered category signal, with a strongest return point of 60.7\%. The full attack union retains both stripped and unstripped traces. This separates the ability to recognize padding phases from the ability to recover a task's operation.
\begin{table}[t]\centering\footnotesize\setlength{\tabcolsep}{3pt}
\caption{Exp 4 intent BA after learned flush removal (95\% customer-cluster interval). Each row: 138 test episodes, 24 customers. Maximum of three classifiers plus majority; the full attack union additionally keeps unstripped observations.}\label{tab:r8-flush}
\begin{tabular}{@{}llrr@{}}\toprule
Prompt & Schedule & Cancel & Return\\\midrule
Standard & Serial & 57.8 [50.0, 70.7] & 53.4 [50.0, 64.6]\\
Standard & Parallel & 50.0 [50.0, 64.0] & 50.0 [50.0, 56.7]\\
Rules & Serial & 50.1 [50.0, 55.6] & 60.7 [50.8, 70.9]\\
Rules & Parallel & 50.8 [50.0, 61.7] & 54.0 [50.0, 67.0]\\
\bottomrule\end{tabular}\end{table}


\paragraph{Match cost to the protected boundary}
Within measured RPC conditions (M), per-call constant framing uses 3.92$\times$ bytes and 16.6\,ms p95, versus serial shadow's 15.11$\times$ and 26.2\,ms; its full-union BA is higher (68.6/62.7\% versus 61.4/53.4\%), a measured cost--inference trade-off. Two-in-flight shadow lowers p95 to 20.8\,ms, versus 56.4/112.3\,ms for 2/5\,ms tunnels. Sessions preserve real responses and final live state; fixed tunnel sizes/counts still expose host-scheduled direction/order. Table~\ref{tab:r7-controls} keeps measured RPC times separate from offline schedule transformations.

Under the common seven attacks, offline random-donor Walkie-Talkie (WT, S) reaches 50.0/56.7\% BA at 2.04$\times$ bytes; measured serial shadow (M) reaches 50.0/52.9\% at 15.11$\times$. This descriptive, shared-metadata comparison identifies a lower-byte shaping candidate. Establishing its online trade-off requires RPC measurements and the extended timing/flush suite; the offline point is neither a theoretical byte lower bound nor demonstrated deployment superiority.

The WF adapters use the same saved traffic and training customers. WT samples donor traces without symmetric site pairing; WTF-PAD samples training inter-arrival histograms; Tamaraw and FRONT supply fixed-rate and front-loaded controls. Input-derived fragments preserve sizes, directions and order, enabling a controlled metadata comparison. Their simulated spans and measured RPC latency remain separate cost objects.

\paragraph{Emulated WAN delay and jitter}
The relay re-runs all 276 test episodes with 20\,ms/5\,ms and 50\,ms/10\,ms one-way delay/jitter. The frozen loopback flush detector transfers with 89.4\% BA (versus 90.4\%), so the cover phase stays recognizable. Jitter does mask micro-timing: serial-shadow full-union cancellation BA falls from 61.4\% to 50.0\% and 54.0\%. This masking depends on observer position; an observer adjacent to the tool server sees loopback-like timing. Fixed horizons fail under delay: the 2\,ms tunnel's 133-slot (266\,ms) horizon completes only 85/276 and 5/276 sessions, whereas shadow cover and per-call framing complete all sessions at 855.4/1,999.5\,ms serial-shadow p95. A tunnel horizon must therefore be sized to path RTT, at a corresponding byte cost.

\textbf{Finding.} The same-capture timing contrast establishes schedule leakage missed by length-only evaluation. Exp~3's 50.0/50.0\% and Exp~4's 61.4/53.4\% also differ in attacker and collection scope. Within measured conditions, cheaper per-call framing leaves higher full-union inference; offline WF adapters identify lower-byte candidates under the common seven. Selecting between transport shaping and dispatch isolation therefore requires both the receiver boundary and the matched cost object.

